\documentclass[11pt]{article}
\usepackage[T1]{fontenc}
\usepackage{lmodern,amsmath,amssymb,amsthm,booktabs}
\usepackage[margin=1in]{geometry}
\usepackage[hidelinks]{hyperref}
\newtheorem{proposition}{Proposition}
\title{A conditional $2^{-59}$ exponent saving for integer multiplication}
\author{Douglas Colkitt\thanks{Prepared with assistance from OpenAI Codex.
This is not independent review or endorsement by OpenAI.}}
\date{October 7, 2026}
\begin{document}
\maketitle
\begin{abstract}
A weighted paired-block circuit reduces the finite bit network, while two
refinements of the downstream accounting permit a larger transform dimension:
the guard bound uses the actual early stopping depth, and the Gaussian width
is chosen to meet its stated interface directly. At ground size $h=50$, these
changes give the conditional saving $\kappa=2^{-59}$, with a strictly positive
exact margin. This is a factor $128/13$ above the preceding $13\cdot2^{-66}$
witness. The underlying multiplication theorem remains an assumption.
\end{abstract}
\section{Scope}
The source patch is against OpenAI's \emph{Integer multiplication below
$n\log n$}, pinned at
\begin{center}\texttt{adc7f1241b42e322a6451854ab7e4b4c146bf78a}.\end{center}
The computational model uses a fixed finite alphabet and a fixed number of
one-dimensional tapes. We retain nonadjacent routing, frames from orthogonal
complements, and auxiliary stage sharing. The ground size is
now 50 for both finite networks. The complex motif uses its original family;
only the bit motif receives the new sum circuit. No practical runtime or
unrestricted optimality claim is made.
\subsection{A paired-block circuit for the bit correction}
\label{sec:paired-bit-construction}
Use $h=50$, $v=\binom h3=19600$, $N=v^3$, and $m=h^3=125000$.
The complex network above is retained at this ground size. For the bit
network, we replace the intersection-one side correction by a smaller circuit.
Fix a common point $i$. On the other $n=h-1=49$ points, inputs $x_{ab}$
represent source triples $\{i,a,b\}$. The required partial output is
\[
 z_{cd}=\sum_{\{a,b\}\cap\{c,d\}=\varnothing}x_{ab}.
\]
All scalar additions are in $\mathbb F_2$, but the computation below always
adds disjoint formal input supports.

It is useful to solve a weighted version. Give a graph an edge value $e_{uv}$
and a vertex value $w_u$, with pairwise disjoint supports. Return its total
value and the values left after deleting each one or two vertices, including
all their incident edges. At the root, take $w_u=0$ and $e_{uv}=x_{uv}$.
Group the ordered vertices consecutively into pairs, with a final singleton
when necessary. A block $A$ becomes a vertex of a smaller graph with weight
\[
 \widetilde w_A=\sum_{u\in A}w_u+\sum_{\{u,v\}\subset A}e_{uv},
 \qquad
 \widetilde e_{AB}=\sum_{u\in A,\ v\in B}e_{uv}.
\]
Recursively compute the smaller graph's total, its value $O_A$ after deleting
block $A$, and its value $F_{AB}$ after deleting both $A,B$.

For $a\in A$, put
\[
 C_a=\sum_{u\in A\setminus\{a\}}w_u,\qquad
 E_{a,B}=\sum_{u\in A\setminus\{a\},\ v\in B}e_{uv},\qquad
 H_{a;B}=C_a+\sum_{C\ne A,B}E_{a,C}.
\]
Compute all $H_{a;B}$ by the leave-one-out prefix/suffix routine on the
vector $(C_a,(E_{a,C})_{C\ne A})$, requesting only the outputs that omit
one of the $E$ terms. The total of this vector is denoted $H_a$.
If $a,b$ are in different blocks $A,B$, their exclusion output is
\[
 F_{AB}+H_{a;B}+H_{b;A}
   +\sum_{u\in A\setminus\{a\},\ v\in B\setminus\{b\}}e_{uv}.
\]
For two vertices in one block it is $O_A$; the single exclusion of $a$ is
$O_A+H_a$. These are disjoint partitions of the surviving weighted graph.
This proves the recursion by induction. The weights created by internal edges
remain inside the recursion, avoiding a separate exclusion circuit for them.

Use direct sums for at most four vertices. Compute other totals by balanced
binary addition. In the four-term formula, first add $F_{AB}+H_{a;B}$,
then add the other two terms, then combine those two results. Intern equal
formal supports and remove unused nodes. At $n=49$, the exact circuit has
9813 additions and 1176 designated outputs. All 1271256 nonzero coefficients,
and every required zero, are checked by comparing all output supports with
their definitions.

Identify equal inputs and equal sums across the $h$ common-point groups.
A sum shared by groups $i\ne j$ consists entirely of triples $\{i,j,k\}$;
its fixed pair and variable vertices identify it exactly. Keep the first
encountered decomposition, in common-point and then local-node order. Retain
all $q=h\binom{h-1}{2}=58800$ partial outputs, injecting their sum into each
physical target. Since every neighbor pair has a unique common point, this
remains exactly the intersection-one matrix. The global graph has
\[
 c=450394,\qquad q=58800,\qquad c+q=509194.
\]
Here 40256 additions have been merged from the separate local graphs.

\subsection{Reversible roles and frames in both directions}
For every node, regard its later uses and designated outputs as outgoing edges.
At an input node, allocate one role per outgoing edge and copy a pivot into
the others. At an addition node, add the second incoming role into the first,
reuse the first as one outgoing pivot, and copy it into fresh roles for the
other outgoing edges. The other input role retires. These are invertible
pointwise gates. There are $2c+q$ outgoing uses and $c$ pivot identifications,
so $R=c+q=509194$ roles suffice. Roles are compiled from the merged graph;
no simultaneously occupied physical slots are identified.

Let $L$ be the product of these gates, $V$ the source copy, $J$ the injection
of the designated partial outputs, and $G,R_0$ the original central operations.
The chronological schedule is
\[
 L,\ J,\ L^{-1},\ R_0,\ V,\ G,\ R_0,\ L,\ J,\ L^{-1},\ G,\ V.
\]
The two injections give $JLz+JL(z+Vx)=JLVx$ for arbitrary initial scratch $z$.
The cleanup restores $z$ and all central scratch. Since $JLV+R_0G=I$,
the invocation is $y\gets y+x$. Reverse and invert the schedule for stage 2.
Forward, inverse, forward invocations therefore exchange the two data banks.

Use the original rational form $I-J/9$ on $F=\mathbb Q^{50}$. At a stage,
write $A=B\mathbin\perp P$ as before, suppressing the future line factor $Q$.
Set $D_0=B\otimes F$, $D_1=A\otimes F$, and
$D_U=(B\otimes F)\mathbin\perp(P\otimes U)$.
Let $U_z$ be the span of the source indicators reaching node $z$.
Every retained node is a sum from an original common-point group, so all its
triples contain some $i$. For a vector $u$ in that span,
\[
 \sum_j u_j=3u_i,\qquad
 \langle u,u\rangle=\sum_{j\ne i}u_j^2>0\quad(u\ne0).
\]
Thus $U_z$ is nondegenerate. Along every graph edge $z\to w$,
$U_z\subset U_w$. Forward middle gates use $D_{U_z}$; the physical roles
follow paths, so these frames increase. A designated output for target $T$
has $U_z\subset t_T^\perp$, because every contributing triple is a neighbor.

For the reversed middle computation use $D_{U_z^\perp}$. The ambient form
is nondegenerate, so all these complements are nondegenerate. Reversing an
edge reverses inclusion: $U_w^\perp\subset U_z^\perp$.
The original output begins at the physical $X$ line
$\langle t_T\rangle\subset U_z^\perp$; an original input ends at exactly
its physical $Y$ complement $t_S^\perp$. No positivity assumption on the
reachable-target spans is needed. At every gate all incident roles receive
the same frame. Early mixers use $D_0$, and late cleanup uses $D_1$.
Newly active roles enter from $D_0$ and retired roles grow to $D_1$.
Tensoring by $Q$ gives the original data-stage boundaries in both orientations.
The only decreases remain the central returns, each of dimension $h$.

Share the complete auxiliary bank of a stage-1 invocation $(A,B)$ with stage 3
at $(B,\pi(A))$, matching local role indices. Here $\pi$ is a bijection of
triples with $|T\cap\pi(T)|=1$: pair consecutive ground points; for a triple
with a full pair and singleton, keep the singleton and cycle the full pair
among pairs not containing it; otherwise keep the point in the least indexed
pair and flip the other two to their partners. Each rule is bijective on its
class. All invocations restore arbitrary scratch, and the joining labels are
\[
 E=F\otimes\langle t_A\rangle\otimes\langle t_B\rangle,
 \qquad H=\langle t_B\otimes t_{\pi(A)}\rangle^\perp\otimes F.
\]
They are nondegenerate and $E\subset H$ by the chosen orthogonality.
Each join removes one auxiliary role and exactly $m$ of edge rank.

\begin{proposition}[Paired-block bit interface]\label{prop:paired-bit-interface}
The new bit network restores arbitrary scratch, exchanges the data banks,
and satisfies the original all-role endpoint identity with
\begin{align*}
 W_{\rm b}^{\rm pair}&=2N+2v^2(R+h)=406321422080000,\\
 s_{\rm b}^{\rm pair}&=W_{\rm b}^{\rm pair}m-N+6v^2h^2
                         =50790175992864000000,\\
 W_{\rm b}^{\rm pair}m-s_{\rm b}^{\rm pair}&=1767136000000>0.
\end{align*}
\end{proposition}
\begin{proof}
The scalar schedule and both frame assignments were verified above. The total
decreasing dimension is $3v^2h^2$. Telescoping gives total absolute label change
$W_{\rm b}^{\rm pair}m-2N+6v^2h^2$; the unchanged negative source projections
add $N$ rational rank. Data endpoints are unchanged and every surviving
auxiliary role has endpoints zero and $I_m$, with its scalar value restored.
The original common-frame identity and rational-shear transfer apply. All
circuits, rational bases, and gates are fixed finite data.
\end{proof}

\subsection{Explicit exponents at ground size fifty}
Take $a=296/10^{11}$, $a_{\rm c}=1/10^{11}$, and $L_0=11737/1000$.
The exact deficits satisfy
\[
 \eta_{\rm b}^{\rm pair}=\frac{23}{661055000}>aL_0,\qquad
 \eta_{\rm c}=\frac{239}{1202215191250}>a_{\rm c}L_0.
\]
For $1\le x\le2$, put $z=(x-1)/(x+1)$,
$S(x)=2\sum_{j=0}^{23}z^{2j+1}/(2j+1)$, and
$E_0(x)=2z^{49}/(49(1-z^2))$. Then
$S(x)\le\log x\le S(x)+E_0(x)$, and exact rational comparison gives
\[
 \log m\le16(S(2)+E_0(2))+S(m/2^{16})+E_0(m/2^{16})<L_0.
\]
The inequality $e^{-x}>1-x$ therefore proves the required strict bounds with
\begin{equation}\label{eq:explicit-motif-exponents}
 \tau=1-\frac{296}{10^{11}},\qquad \sigma=1-\frac1{10^{11}}.
\end{equation}
The complex network itself is unchanged apart from the ground size $h=50$;
its construction, binary residual bases, and scalar corrections remain those
proved above.

\section{The guard refinement}
\subsection{A guard bound using the actual stopping depth}
\label{sec:stopped-guard}
The complex network is still evaluated exactly until a complete normalized
butterfly layer is finished. Write $W=W_{\rm c}$ and $s=s_{\rm c}$ for that
network, and retain
\[
 E=64(W+m+1)^3,\quad B=s+E,\quad C_0=128mB^2.
\]
At $h=50$ the exact values are
\[
 m=125000,\quad W=369474390296480000,\quad
 s=46184298777878576000000,
\]
so $m\ge3$ and $2\le s<m^5$ by integer comparison. We will use $C_1=2$ and
$\Delta=\lceil C_0d^2\rceil$ for every rational $9/10\le\beta<1$.

Measure dependency depth by complex additions, subtractions, halvings, sign
changes, and multiplication by $i$. A chain follows an operand into a later
result. Concatenating every scalar gate and recursive edge call is a valid
upper bound, even though most do not belong to one dependency chain.
In the retained complex motif each wire touches at most twelve scalar gates,
so there are at most $12W$ gates. A gate output sums at most $W$ terms,
with at most a halving and an addition per term. Evaluating all outputs from
saved inputs uses at most $2W^2$ operations per gate. Inverse-child corrections
cost at most $4s$, and source, sink, and bank corrections at most $4W+4$.
Thus their total is bounded by $24W^3+4s+4W+4<E$.
The new bit motif performs exact permutations of coefficient encodings; it
does not introduce coefficient arithmetic into this depth count.

The individual $C$-kernel formula uses at most eight elementary operations.
Consequently the same conservative recurrence remains valid:
\[
 A(e)\le sA(e/m)+E\quad\hbox{at an internal node},\qquad
 A(e)\le8e\quad(e<d^\beta).
\]
For a root $e\le d$, the actual number $j$ of internal levels is at most
$(1-\beta)\log_m d+1$. If $j>0$, this follows from the last internal size
$e/m^{j-1}\ge d^\beta$; for $j=0$ it is immediate. Hence
\[
 s^j\le s\,d^{5(1-\beta)},\qquad
 A(e)\le8d^\beta s^j+E\sum_{i<j}s^i
       \le s(8+E)d^{5-4\beta}
       \le9B^2d^{7/5}.
\]
Here $s\ge2$, $d\ge1$, $\beta\ge9/10$, and
$s(8+E)\le9B^2$; the estimate also covers a root that is already a leaf.
Keeping $E$ additive is essential: it is a fixed cost at a node, not another
branching factor.

The base-$m$ partition has at most $(m-1)(1+\lfloor\log_m d\rfloor)$
pieces. Since $m\ge3$ and $\log d\le\sqrt d$ for $d\ge1$, this is at most
$2m\sqrt d$. Concatenating their depths therefore costs at most
$18mB^2d^{19/10}$. Preprocessing uses at most $d$ individual kernels, costing
$8d$. The two outer diagonals and the scalar $b^D$ contribute at most
$t+9$ further units, where $t=\lfloor D/2\rfloor$; a shift by $t$ is charged
$t$ denominator bits. Together $8d+t+9\le18d$. Thus the total depth and
shift charge are bounded by
\[
 18mB^2d^{19/10}+18d\le36mB^2d^2<C_0d^2.
\]
As in the original induction, an input in $2^{-p}\mathbb Z[i]$ with modulus
at most one has denominator exponent at most $p+\Delta$ and magnitude at most
$2^\Delta$ throughout. Reserving $p+\Delta$ fractional bits and $\Delta+3$
integer and sign bits per real component therefore represents all values
exactly. Address-swap routines restore complete encodings before coefficient
arithmetic resumes and do not change this argument. If $2\epsilon<1$, then
$\Delta=o(p)$ and the full coefficient width remains $O(p)$.

\section{A smaller Gaussian width}
Keep $p=6b$, $\eta=1/(4d)$, and the same prime intervals. They still give
$\theta_i=t_i/s_i-1>1/(4d)$. Replace the earlier Gaussian width by
\[
 \alpha=\left\lceil(32db)^{1/4}\right\rceil,\qquad
 \gamma=2d\alpha^2.
\]
Then $\alpha^4\theta_i>8b>p$, precisely the strict hypothesis needed by
the unchanged one-dimensional resampling interface. Also
\[
 \alpha\le2(32db)^{1/4},\qquad
 \gamma<46d^{3/2}\sqrt b.
\]
For $d=\Theta(p^\epsilon)$, we now have
$\alpha=\Theta(p^{1/4+\epsilon/4})$ and
$\gamma=O(p^{1/2+3\epsilon/2})$.
Thus $\epsilon<1/3$ implies $\gamma=o(p)$ and
$2\le\alpha<\sqrt p$ eventually. The previous convenient condition
$\epsilon<1/12$ and its bound $\gamma\le28b^{2/3}$ are not needed.

The unchanged Gaussian line-map cost $O(dTp^{3/2+\delta}\alpha)$ now has
normalized exponent $3/4+\delta+5\epsilon/4$. The other six assembly
exponents are unchanged. Prime selection still requires $1-2\epsilon>0$;
the product of the selected primes remains between $T/2$ and $T$, with the
same separation and coprimality guarantees. The new width satisfies every
stated resampling hypothesis without changing that analytic lemma.

For the chosen $\epsilon=199/1000$, the exponents of $\alpha$ and $\gamma$
are respectively $1199/4000$ and $1597/2000$. The explicit cutoff
$b\ge2^{40}$ ensures $46b^{1597/2000}\le b/4$ by integer comparison
$184^{2000}<2^{16120}$. Replacing the former $2^{24}$ cutoff by $2^{40}$
therefore preserves the source-scaling and final rounding bounds, together
with the other existing eventual cutoffs. All grid, contraction, and error
lemmas are retained with their original hypotheses verified.

\section{Exact parameters and final margin}
With $a=296/10^{11}$ and $a_{\rm c}=1/10^{11}$ take
\begin{gather*}
 \tau=1-a,\quad\sigma=1-a_{\rm c},\quad\beta=999/1000,
 \quad\epsilon=199/1000,\quad\delta=1/10000,\quad C_1=2,\\
 c=999a/1000,\quad\lambda=1-1999a^2/2000,
 \quad\lambda'=1-999a^2/1000.
\end{gather*}
Here $\tau(1+c/\beta)=1-a^2<\lambda<\lambda'<1$;
$\sigma<\lambda$ and $\sigma+\beta(1-\sigma)<\lambda'$ hold by exact
rational comparison. The stopping test is $e^{1000}<d^{999}$.
The seven margins, with nonadjacent routing and the new Gaussian width, are
\begin{align*}
 g_1&=1-\epsilon(1+c),&g_2&=\epsilon c(1-\tau),\\
 g_3&=\epsilon(1-\lambda'),&g_4&=(1-\tau)(1-\epsilon),\\
 g_5&=1/4-\delta-5\epsilon/4,&g_6&=1-\delta-\epsilon,\qquad g_7=\epsilon.
\end{align*}
Exact arithmetic gives
\[
 G=\min_j g_j=g_2=g_3
 =\frac{272158569}{156250000000000000000000000}>2^{-59}.
\]
The ratio $G/2^{-59}$ is approximately $1.00408789$; the displayed strict
rational comparison, not this decimal, is the certificate. The fixed gap
absorbs the remaining fixed powers of $\log p$.

The guard exponent is $2\epsilon=199/500<1$, the prime-interval growth
exponent is $301/500>0$, and $\ell=\Theta(p^{801/1000})$.
We have $K=\Theta(p^{\epsilon c})=o(\ell)$ and $K/\log p\to\infty$.
The Gaussian margin is $23/20000>0$. Thus all setup, stopping, precision,
layout, and absorption conditions in the patched proof hold, yielding
\[
 T(n)=O\bigl(n(\log n)^{1-2^{-59}}\bigr)
\]
under the retained upstream assumptions.

\section{Reproducibility and boundary of verification}
The reproducible artifacts are
\begin{itemize}
\item \path{certificates/paired-network.json};
\item \path{patches/h50-paired-59.patch};
\item \path{scripts/paired_exclusion_circuit.py};
\item \path{scripts/shared_point_circuit.py};
\item \path{scripts/paired_network.py}.
\end{itemize}

The checks compare every local and global output support exactly, verify both
frame directions on every compiled role, and check small forward and reverse
invocations on every input basis vector, including dirty scratch. Full small
three-stage models check bank exchange and shared auxiliary restoration.
Separate tests cover weighted recursion, stopping-depth bounds, guard-model
selection, revised assembly constraints, and source-patch dependencies.
These checks support the written proofs; they do not independently establish
or formally verify the complete upstream algorithmic theorem.
\end{document}
